\documentclass[10pt,a4paper]{amsart}
\usepackage[margin=19mm]{geometry}
\usepackage{amsmath,amssymb,mathtools}
\usepackage[scr=rsfs,cal=euler]{mathalfa}
\usepackage{xfrac}
\usepackage[unicode,hidelinks,bookmarks=false]{hyperref}
\newcommand{\ie}{i.e.\ }
\newcommand{\cf}{cf.\ }
\newcommand{\resp}{respectively\ }
\newcommand{\Subsection}{\subsection}
\providecommand{\nobreakdash}{\nobreak\hbox{-}}
\setlength{\emergencystretch}{2em}
\multlinegap0pt
\usepackage{bm}
\usepackage{bbm}
\usepackage{dsfont}
\usepackage{xspace}
\usepackage{cancel}
\usepackage{mathtools}
\usepackage{amsthm}
\makeatletter
\@ifundefined{theorem}{\newtheorem{theorem}{Theorem}}{}
\@ifundefined{proposition}{\newtheorem{proposition}[theorem]{Proposition}}{}
\makeatother
\title[Conjectural Decidability of the Skolem Problem]{Conjectural Decidability of the Skolem Problem}
\author{Florian Luca, Jo\"el Ouaknine, James Worrell}
\date{Source: arXiv:2607.15510v1 (CC BY 4.0)}
\begin{document}
\maketitle

\noindent\emph{Source and licence.} Conjectural Decidability of the Skolem Problem. Version arXiv:2607.15510v1 (CC BY 4.0), distributed under the Creative Commons Attribution 4.0 International licence, as recorded in the arXiv licence metadata for this exact version and shown on the abstract page https://arxiv.org/abs/2607.15510v1. Full rights notice: licenses/arxiv-2607.15510-CC-BY-4.0.txt.\par\smallskip

\noindent\textit{Editorial scope.} Counted unit: the Proposition whose source label is \texttt{prop:density} in the article (source0.tex lines 574--581, source label \texttt{prop:density}), delivered together with the article's own complete original proof of it (source0.tex lines 583--650, one \texttt{proof} environment of 68 lines). No mathematics is written, repaired or re-derived: statement and proof are the article's own text line for line. Required context retained uncounted: eq:house (lines 535--539); eq:21 (lines 547--550). Every definition, display and label of those lines is retained unchanged. Omitted from this package and not redistributed: 1--534, 540--546, 551--573, 582--582, 651--1018. Nothing inside a retained excerpt is omitted or altered: every retained body is delivered exactly as the article's lines are written, so no omission receipt of any kind accompanies this package. Citations: the retained excerpts carry no citation command at all, so no bibliography entry of the article is transcribed and no reference is redistributed. Reference adaptation: none was needed and none was made; every reference inside the delivered unit resolves inside it, so no unresolved reference or citation is produced. Printed slips kept literal: no printed word, symbol, hypothesis or number of the counted statement or of its proof was altered. Layout and class: the article's own class is \texttt{elsarticle} with packages including amsmath, amssymb, amsfonts, algorithmic, graphicx, textcomp, xcolor, amsthm, tikz, hyperref; the wrapper is the lane's pinned \texttt{amsart} editorial wrapper at 10pt on A4, which restates the theorem-like environments the retained text needs and adds the article's own macro definitions that the retained text needs (none), with extra wrapper packages (bm, bbm, dsfont, xspace, cancel, mathtools). The delivered numbering is the unit's own. Assets: no retained span contains a figure environment, a graphics-inclusion command, a PDF-page inclusion, a table or a TikZ picture; no image file is copied into this package, the package declares no asset dependency and no image file is redistributed with it. Licence: version arXiv:2607.15510v1 of this article is distributed under the Creative Commons Attribution 4.0 International licence, as recorded in the arXiv licence metadata for this exact version and shown on the abstract page https://arxiv.org/abs/2607.15510v1. A full rights notice is delivered at licenses/arxiv-2607.15510-CC-BY-4.0.txt. Characters: every retained excerpt is pure ASCII, so the delivered engine, XeLaTeX with its default Latin Modern text font, reports no missing character.
\begin{equation}
\label{eq:house}
  |\Delta| (k^k(1+H)^k)^k < (k(1+H))^{2k^2} < (1+H)^{4H^2}< H^{H^3}
  \, .
\end{equation}  

\begin{equation}
\label{eq:21}
v_{m,\sigma}=\sum_{i=1}^s P_i(m)\beta_i \alpha_i^{m} \, .
\end{equation}

\begin{proposition}
\label{prop:density}  
We have 
\[
\# {\mathcal P}_{\text{\rm bad}}(X)<X^{2/3}
\]
for all $X>X_0$, where $X_0$ is some effective absolute constant. 
\end{proposition}

\begin{proof}
In order to estimate the size of ${\mathcal P}_{\text{\rm bad}}(X)$,
we first need to find out: 
\begin{enumerate}
\item[\textbf{1}.] How many such expressions \eqref{eq:21} are there?
\item[\textbf{2}.] How large are they?
\end{enumerate}

For (\textbf{1}), let us count the number of distinct possible LRS of
size at most $H$. Such LRS have coefficients $a_1,\ldots,a_k$ and initial values
$u_0,\ldots,u_{k-1}$ all in $[-H,H]$, an interval containing at
most $2H+1<3H$ integers. Altogether for fixed $k$ there are at most
$(3H)^{2k}\le (3H)^{2H}$ $2k$-tuples, and summing up over $k$ we
derive an upper bound of $H(3H)^{2H}<H^{3H}$ distinct
possible LRS of size at most $H$.

This in turn is an upper bound on
the number of $s$-tuples $((Q_i ,\alpha_i))_{i=1}^s$. We must then
multiply this quantity with the number of possible permutations of the
characteristic roots,
%$(\beta_1,\ldots,\beta_s)$,
which is at most $H!<H^H$. There are therefore
at most $H^{4H}$ linear recurrence sequences
$\mathbf{w} = \langle w_m \rangle_{m=0}^{\infty}$ whose $m$-th term is
given by
\[
w_m=\sum_{i=1}^s P_i(m)\beta_i\alpha_i^m\quad {\text{\rm
    for~all}}\quad m\ge 0 \, .
\]

This answers (\textbf{1}).
As for (\textbf{2}), recall that the coefficients of $P_i$ are of
absolute value at most 
$H^{H^3}$, as per~\eqref{eq:house}. $P_i(m)$ comprises at most $H$ monomials, the largest of which
is at most $m^H < X^H$, and the largest root has
magnitude at most
$1+H < 2H$.
Thus each individual term $w_m$ is of absolute value at most
\begin{multline*}
  H^{H^3+1} (2H) X^H (2H)^{X^{1/4}} = \\
  \exp\left((H^3+1)\log H + \log (2H) + H \log X + X^{1/4}\log (2H)
  \right)\\ <\exp(X^{0.26})
\end{multline*}
for $X>X_0$, since $H$ is tiny in comparison to $X$.
Hence the norm of the number shown in \eqref{eq:21} is of size at most 
\[
\exp(H! X^{0.26})<\exp(X^{0.27})\quad {\text{\rm for}}\quad X>X_0 \,
,
\]
since the degree of $\mathbb{K}$ is at most $H!$ (as $\mathbb{K}$ is
the splitting field of a polynomial of degree at most $H$).
Moreover, as noted
earlier, there are at most $H^{4H}$ such expressions.
Thus a bad prime $P$ divides an integer which is a product of such numbers and is of size at most 
\[
\exp(H^{4H} X^{0.27})<\exp(X^{0.28})\quad {\text{\rm for}}\quad
X>X_0 \, .
\]

Therefore the number of possible choices for $P$ is at most $X^{0.28}$. Since
the number of choices for $m$ is at most $X^{0.25}$, we conclude that,
for $X > X_0$,
the cardinality of $\mathcal{P}_{\text{\rm bad}}(X)$ is at most
\[
X^{0.25+0.28} < X^{2/3} \, ,
\]
as required.
\end{proof}

\end{document}
